\section{Consequences and boundaries}
\begin{corollary}\label{cor:basis13}
Under the analytic inputs of Section~\ref{sec:inputs}, the set
\(C_{12}\cup\{0\}\) is an asymptotic basis of exact order
thirteen: every sufficiently large integer is a sum of at most
thirteen elements of \(C_{12}\), and twelve do not suffice.
\end{corollary}
\begin{proof}
For a large target choose \(j\in\{0,\ldots,11\}\) with
\(N-j\equiv2\pmod{12}\). Theorem~\ref{thm:c12} writes \(N-j\)
as two core values; add \(j\) copies of one. This uses at most thirteen
positive summands. Conversely, \(25+60r\), \(r\ge0\), is
\(1\pmod{12}\) and divisible by the forbidden prime five.
At most twelve positive core summands would force exactly one
summand on that residue class, which is impossible. Padding with
zeros gives the equivalent fixed-order assertion.
\end{proof}

The field formulation follows from classical abelian class field
theory. A subgroup condition on Frobenius is complete splitting in
its fixed field. Finite conjunctions become intersections of inverse
image subgroups in the compositum. Kronecker--Weber places this
compositum in a cyclotomic field
\cite[Chapter I, Theorem 4.16]{milnecft}, giving precisely the
fixed subgroups in Theorem~\ref{thm:general}, after discarding the
fixed bad-prime set. The reduction is classical; the pointwise
indicator correlation is the analytic assertion used here.

Exact classifications and their finite certificates are treated in
the companion \emph{Exact Sumsets of Primitive Quadratic Norms}
\cite{exactcompanion}.
The separate article
\href{https://hexagonmath.org/2611.00593v1}{\emph{Pointwise representation
counts for Eisenstein restricted sums}} concerns the placement of summands.
Higher-fold sums and independent coupled systems can use classical
linear-form correlations, with a different proof from this reflected
binary problem. Growing-shift uniformity and higher indicator correlations
are not covered by the statements here.
